% =====================================================================
%  Paper 1, Appendix: Six-Degree-of-Freedom Equations of Motion
%  Owns: apx:eom
%
%  The vehicle is a rigid body with attitude and body rates as STATES, 
%  and (alpha, beta) as OUTPUTS of the attitude solution rather than 
%  commanded quantities. Controls are flap deflections and RCS moments.
%
%  Why not 3-DOF: a point-mass model treats (alpha, sigma) as
%  instantaneously commandable, which assumes an attitude loop that
%  tracks perfectly and a trim map that is always solvable. Neither is
%  true of a real vehicle.
% =====================================================================

\section{Six-Degree-of-Freedom Equations of Motion}
\label{apx:eom}

This appendix records the rigid-body dynamics the body of the paper
treats as given. Nothing here is a contribution; the derivation is
standard entry mechanics\todo{cite sources} assembled once with a fixed set of 
frames and a fixed parametrization. Frame errors are the classic $6$-DOF defect 
and are invisible in the assembled algebra, so the transformations are stated 
explicitly and used without exception.

\subsection{State vector and frames}
\label{sub:apx_state_and_frames}

\paragraph{Frames.} Four frames are used throughout.
\begin{itemize}
  \item $\mathcal I$: Earth-centred inertial (ECI), non-rotating.
  \item $\mathcal E$: Earth-centred Earth-fixed (ECEF), rotating with the
        constant planetary rate $\vec\omega_E = \omega_E\,\hat z_{\mathcal I}$.
  \item $\mathcal L$: local geocentric frame at the vehicle, with unit
        vectors $(\hat e_r,\hat e_\theta,\hat e_\phi)$ pointing radially
        outward, east, and north.
  \item $\mathcal B$: body frame, fixed at the instantaneous centre of
        mass, axes along the vehicle's principal design directions.
\end{itemize}
The attitude quaternion $\vec q$ carries $\mathcal I \to \mathcal B$; the
geocentric angles $(\theta,\phi)$ carry $\mathcal I \to \mathcal L$.

The
direction-cosine matrix built from $\vec q$,
\begin{equation}
\label{eq:apx_dcm}
\mathbf R_{B/I}(\vec q)=
\begin{pmatrix}
q_0^2+q_1^2-q_2^2-q_3^2 & 2(q_1q_2+q_0q_3) & 2(q_1q_3-q_0q_2)\\
2(q_1q_2-q_0q_3) & q_0^2-q_1^2+q_2^2-q_3^2 & 2(q_2q_3+q_0q_1)\\
2(q_1q_3+q_0q_2) & 2(q_2q_3-q_0q_1) & q_0^2-q_1^2-q_2^2+q_3^2
\end{pmatrix},
\end{equation}
is polynomial of degree two in $\vec q$, which is the property
\cref{apx:embedding} exploits. The geocentric map is
\begin{equation}
\label{eq:apx_RLI}
\mathbf R_{L/I}(\theta,\phi)=
\begin{pmatrix}
\cos\phi\cos\theta & \cos\phi\sin\theta & \sin\phi\\
-\sin\theta & \cos\theta & 0\\
-\sin\phi\cos\theta & -\sin\phi\sin\theta & \cos\phi
\end{pmatrix}.
\end{equation}

\paragraph{State vector.} Before augmentation the state carries thirteen
components,
\begin{equation}
\label{eq:apx_state}
x=\bigl(\underbrace{r,\theta,\phi}_{\text{position}},\;
        \underbrace{u,v,w}_{\text{body velocity}},\;
        \underbrace{q_0,q_1,q_2,q_3}_{\text{attitude}},\;
        \underbrace{p,q,r_{\!b}}_{\text{body rates}}\bigr)\in\mathbb R^{13},
\end{equation}
augmented by mass $m$ when a propulsive mode is present
(\cref{sub:apx_per_mode_forms}), and by the recession depth $s$ and modal
state $(\eta,\dot\eta)$ once the aerothermoelastic loop is closed
(\cref{sub:apx_rotational,sub:certified_envelopes}). Here $(r,\theta,\phi)$
are geocentric radius, longitude and latitude; $(u,v,w)$ are the body-axis
components of the \emph{planet-relative} velocity
\begin{equation}
\label{eq:apx_relvel}
\vec V \;=\; \vec V_{\mathcal I} - \vec\omega_E\times\vec r - \vec w_{\mathrm{wind}},
\end{equation}
with $\vec V_{\mathcal I}$ the inertial velocity and $\vec w_{\mathrm{wind}}$
the local wind (\cref{sub:apx_atmosphere}); and $(p,q,r_{\!b})$ are the body
components of the inertial angular velocity $\vec\omega\equiv\vec\omega_{B/I}$.

\subsection{Translational equations}
\label{sub:apx_translational}

Newton's law holds in $\mathcal I$: $m\,\dot{\vec V}_{\mathcal I}\big|_{\mathcal I}=\vec F$,
with $\vec F=\vec F_g+\vec F_a+\vec F_t$ (gravity, aerodynamic, thrust).
We want the rate of the \emph{relative} velocity resolved in $\mathcal B$.
Writing $\vec V_{\mathcal I}=\vec V+\vec\omega_E\times\vec r$ and
differentiating in $\mathcal I$,
\begin{equation}
\dot{\vec V}_{\mathcal I}\big|_{\mathcal I}
=\dot{\vec V}\big|_{\mathcal I}+\vec\omega_E\times\vec V_{\mathcal I}
=\underbrace{\dot{\vec V}\big|_{\mathcal B}+\vec\omega\times\vec V}_{\text{transport of }\vec V}
+\vec\omega_E\times\vec V+\vec\omega_E\times(\vec\omega_E\times\vec r),
\end{equation}
where the transport theorem $\tfrac{d}{dt}\big|_{\mathcal I}=\tfrac{d}{dt}\big|_{\mathcal B}+\vec\omega\times$
was applied to $\vec V$. Solving for the body-frame rate,
\begin{equation}
\label{eq:apx_trans_vector}
\boxed{\;
\dot{\vec V}\big|_{\mathcal B}
=\frac{1}{m}\bigl(\vec F_a+\vec F_t\bigr)+\vec g
-\bigl(\vec\omega+\vec\omega_E\bigr)\times\vec V
-\vec\omega_E\times(\vec\omega_E\times\vec r)\; }
\end{equation}
with $\vec F_g=m\vec g$. The term $(\vec\omega+\vec\omega_E)\times\vec V$
collects body-rotation transport and Coriolis; expanding
$\vec\omega=\vec\omega_{B/E}+\vec\omega_E$ recovers the familiar
$\vec\omega_{B/E}\times\vec V+2\,\vec\omega_E\times\vec V$ split, the second
being the Coriolis acceleration proper.

\subsection{Rotational equations}
\label{sub:apx_rotational}

Angular momentum about the center
of mass is $\vec h=\mathbf I\vec\omega$, and $\dot{\vec h}\big|_{\mathcal I}=\vec M$;
transporting to $\mathcal B$,
\begin{equation}
\label{eq:apx_euler}
\boxed{\;
\mathbf I\,\dot{\vec\omega}
=\vec M-\vec\omega\times(\mathbf I\vec\omega)-\dot{\mathbf I}\,\vec\omega\;}
\end{equation}
with $\vec M=\vec M_a+\vec M_{\mathrm{rcs}}$. The gyroscopic term is
quadratic in $\vec\omega$, hence polynomial once $\mathbf I$ is; the
inversion $\dot{\vec\omega}=\mathbf I^{-1}(\cdots)$ is carried as the linear
solve $\mathbf I\dot{\vec\omega}=(\cdots)$ so that no rational function of the
state enters the embedding.

\begin{remark}[The inertia tensor is not constant]
Ablation removes mass from the nose and structural deformation
redistributes it, so $\mathbf I=\mathbf I(s,\eta)$ through the recession depth
$s$ and modal amplitudes $\eta$ (\cref{sub:certified_envelopes}), and
Euler's equations acquire the $\dot{\mathbf I}\vec\omega$ term in
\cref{eq:apx_euler}.
We carry it rather than drop it: the whole point of
going to $6$-DOF is to keep the deformation$\to$incidence$\to$moment$\to$attitude
loop closed, and that loop runs through $\mathbf I(s,\eta)$. Its magnitude is
bounded by the recession and modal-amplitude envelopes, so carrying it costs
one Jacobian block, not a new uncertainty.
\end{remark}

\subsection{Aerodynamic forces and moments}
\label{sub:apx_aero}

\subsection{Atmosphere}
\label{sub:apx_atmosphere}

\subsection{Per-mode forms}
\label{sub:apx_per_mode_forms}

\subsection{Boundary discipline}
\label{sub:apx_boundary_discipline}
